% Section 8: Relation to the literature (from research/G1/RH/spectral/NOVELTY_JZC.md; the bibliographic data of the entries
% added for this section were relayed second-hand and must be checked against the sources before submission; DixitKumar:
% statement quoted in research/G1/RH/CITATIONS.md (Hurwitz-DK); metadata from the arXiv abstract page via WebFetch, 2026-10-22).
% Rule R review B of JZC_M1 (B6, 2026-10-23): Rakhmanov, DragnevSaff and VanAsscheVanBaelen are quoted with hypotheses in
% research/G2/CITATIONS.md, entries (R96), (DS97), (VAVB); BeckermannKuijlaars2002 is not quoted in any CITATIONS.md (hence
% "according to"); BeckermannKuijlaars2001 has not been read.
\section{Relation to the literature}\label{sec:literature}

\emph{Constrained equilibrium.} Existence, uniqueness and the structure (band, saturated region, balayage) of equilibrium
problems with an upper constraint are classical \cite{Rakhmanov,DragnevSaff}; these results assume a finite constraint
(Dragnev--Saff) or a compact interval (Rakhmanov) and are not used in the proof of Theorem~\ref{thm:Q}, which establishes
existence and uniqueness directly. For conjugate-gradient convergence, according to \cite{BeckermannKuijlaars2002},
Beckermann and Kuijlaars derived endpoint equations of the same Cauchy-integral form as~\eqref{eq:Q}, with an external
field and compactly supported constraints; a formula for the case without external field is
attributed to \cite{BeckermannKuijlaars2001}, which we have not been able to consult. We have not found the explicit density
formula of Theorem~\ref{thm:Q}, the criterion that $y\sigma(y)$ be decreasing, or the treatment of a constraint of infinite
mass with a hard edge at $0$. The fully saturated equilibrium of \cite{VanAsscheVanBaelen} ($y\sigma$ constant on a
$q$-lattice) is the boundary case of hypothesis~(H2).

\emph{Spectral data and Jacobi parameters.} Orthogonal polynomials on reciprocal spectra, with asymptotics at fixed argument
expressed through a regularised characteristic function, are studied in \cite{StovicekSwiderski}. We know of no expansion
of the recurrence coefficients themselves in terms of spectral zeta values at negative arguments, as in Theorem~\ref{thm:K},
and no reading of geometric spectral invariants (perimeter, corners) from Jacobi parameters. The Lommel family
\cite{VA} is the classical exactly solvable case (Proposition~\ref{prop:bessel}).

\emph{Expansions in a shifted index.} For varying weights $e^{-nV}$ the recurrence coefficients have expansions in even
powers of $1/n$ after a half-index shift; the odd terms are removed by a reindexing symmetry of the string equations
\cite{BleherIts}, and for even weights the coefficients are generating functions of maps \cite{EMP}. For Jacobi-type weights
on $[-1,1]$ the first correction vanishes \cite{KMVV}. The present setting is a fixed discrete measure without a string
equation: the shift involves the spectral constant $Z(0)$, and the evenness comes with the zeta layer and its own even
dressing (Theorem~\ref{thm:N6}, Conjecture~\ref{conj:S1Nprime}). Theorem~\ref{thm:reflect} plays the role of the
reindexing symmetry of \cite{BleherIts}, and the evenness to all orders (Corollary~\ref{cor:even}) is obtained without a string
equation, from a second steepest-descent frame at a non-integer index (Lemma~\ref{lem:frame}). We have not searched the literature
systematically for this device, which is a gauge $z^{\gamma\sigma_3}$ of the exact Riemann--Hilbert problem combined with a change of
frame. The centring of $b_n$ at $N$ and of $a_n$ at $N+\frac12$ (Conjecture~\ref{conj:diag}) is the arrangement of
\cite[Theorem~5.2]{BleherIts}. The identities of the Toda hierarchy used in Proposition~\ref{prop:toda} are elementary and proved in
the text.

\emph{Hurwitz zeta values.} The band-factor moments of a Hurwitz lattice come out of the analysis as Fourier series, and their
identification with $\zeta(-jp,a)$ is Hurwitz's formula, which we quote from \cite[Theorem~1.1]{DixitKumar}; there it is attributed to
Hurwitz and quoted from Titchmarsh's monograph, which we have not consulted. The index transfer of Section~\ref{sec:hurwitz-thm}, a
change of the Szeg\H{o} exponent at fixed Bessel order, and the statements of Theorems~\ref{thm:H} and~\ref{thm:kappa}, have not been
checked against the literature on discrete orthogonal polynomials beyond the works cited in this section.
